% Deutsche Parallelfassung zu geometry/curvature_classification.tex

\chapter{Krümmungsklassifikation}
\label{chap:curvature-classification}

\section*{Motivation}

Die klassische Theorie des Eisenbahn-Alignments klassifiziert Übergänge
üblicherweise nach ihren analytischen Formeln oder ihrer historischen
Ingenieurverwendung. Innerhalb von AIM werden Übergangs- und
Alignment-Strukturen intrinsisch im Krümmungsraum klassifiziert.

Dies ermöglicht eine Klassifikation nach Stetigkeit, Krümmungsverhalten,
Dynamik, Realisierungsverhalten, spektraler Struktur, Laufzeitrobustheit und
betrieblicher Eignung. Krümmungsklassifikation wird damit zu einem
strukturellen Interpretations-Framework statt nur zu einem Katalog von
Kurvenformeln.

\begin{principle}[Klassifikation als Interpretation]
\label{princ:classification-as-interpretation}
Innerhalb von AIM ist Klassifikation kein Benennungsschema für Kurvenformeln,
sondern eine Interpretationsschicht für Krümmungsverhalten.
\end{principle}

\section{Klassifikationsprinzip}

\begin{principle}[Intrinsisches Klassifikationsprinzip]
\label{princ:intrinsic-classification}
Innerhalb von AIM werden Alignments nach ihrem intrinsischen
Krümmungsverhalten und nicht nach rekonstruierter Positionsgeometrie
klassifiziert.
\end{principle}

Zwei geometrisch ähnliche Kurven im Positionsraum können sehr verschiedenes
Krümmungsverhalten besitzen. Umgekehrt können geometrisch unterschiedliche
Realisierungen derselben Krümmungsklasse angehören, wenn sich ihre
intrinsischen Krümmungsverläufe gleichwertig verhalten. Die Klassifikation
erfolgt daher im Krümmungsraum, bevor sie im Positionsraum interpretiert wird.

\section{Stetigkeitsklassifikation}

\subsection{\texorpdfstring{$G^0$}{G0}-Geometrie}
\begin{definition}[\(G^0\)-stetige Geometrie]
\label{def:g0-geometry}
Eine Geometrie ist \(G^0\)-stetig, wenn ihre Position stetig ist.
\end{definition}

\subsection{\texorpdfstring{$G^1$}{G1}-Geometrie}
\begin{definition}[\(G^1\)-stetige Geometrie]
\label{def:g1-geometry}
Eine Geometrie ist \(G^1\)-stetig, wenn ihre Tangentenrichtung stetig ist.
\end{definition}

\subsection{\texorpdfstring{$G^2$}{G2}-Geometrie}
\begin{definition}[\(G^2\)-stetige Geometrie]
\label{def:g2-geometry}
Eine Geometrie ist \(G^2\)-stetig, wenn ihre Krümmung stetig ist.
\end{definition}
Bei der Modellierung von Eisenbahn-Alignments ist \(G^2\)-Stetigkeit häufig
die Mindestanforderung für dynamisch akzeptable Übergänge.

\subsection{\texorpdfstring{$G^3$}{G3}-Geometrie}
\begin{definition}[\(G^3\)-stetige Geometrie]
\label{def:g3-geometry}
Eine Geometrie ist \(G^3\)-stetig, wenn ihr Krümmungsgradient stetig ist.
\end{definition}

Höhere Stetigkeitsklassen verbessern Ruckverhalten, Realisierungsglätte,
Laufzeitstabilität und Fahrzeugreaktion. Die Stetigkeitsklassifikation
verbindet klassische Geometrienotation mit der Krümmungsrauminterpretation.

\section{Klassifikation des Krümmungsverhaltens}

\subsection{Klasse konstanter Krümmung}
\begin{definition}[Klasse konstanter Krümmung]
\label{def:constant-curvature-classification}
Ein Krümmungsgesetz gehört zur konstanten Klasse, wenn
\[
\kappa(s)=\kappa_0.
\]
\end{definition}

\subsection{Lineare Krümmungsklasse}
\begin{definition}[Lineare Krümmungsklasse]
\label{def:linear-curvature-class}
Ein Krümmungsgesetz gehört zur linearen Klasse, wenn
\[
\kappa(s)=a s+b.
\]
\end{definition}
Die klassische Klothoide gehört zu dieser Klasse.

\subsection{Polynomiale Krümmungsklasse}
\begin{definition}[Polynomiale Krümmungsklasse]
\label{def:polynomial-curvature-class}
Ein Krümmungsgesetz gehört zur polynomialen Klasse, wenn
\[
\kappa(s)=\sum_i a_i s^i.
\]
\end{definition}

\subsection{Trigonometrische Krümmungsklasse}
\begin{definition}[Trigonometrische Krümmungsklasse]
\label{def:trigonometric-curvature-class}
Ein Krümmungsgesetz gehört zur trigonometrischen Klasse, wenn es sinus- oder
kosinusbasierte Anteile enthält.
\end{definition}

\subsection{Lookup-basierte Krümmungsklasse}
\begin{definition}[Lookup-basierte Krümmungsklasse]
\label{def:lookup-curvature-class}
Ein Krümmungsgesetz gehört zur Lookup-basierten Klasse, wenn es durch
abgetastete oder tabellierte Krümmungswerte sowie Interpolations- und
Rekonstruktionsregeln definiert ist.
\end{definition}

\subsection{Hybride Krümmungsklasse}
\begin{definition}[Hybride Krümmungsklasse]
\label{def:hybrid-curvature-class}
Ein hybrides Krümmungsgesetz verbindet parametrische und Lookup-Struktur,
Messdaten und abgetastete Korrekturen.
\end{definition}

\section{Symmetrieklassifikation}

\subsection{Symmetrische Übergänge}
\begin{definition}[Symmetrischer Übergang]
\label{def:classification-symmetric-transition}
Ein normierter Übergang ist symmetrisch, wenn sein Krümmungsverlauf bezüglich
der Mitte des normierten Intervalls symmetrisch ist.
\end{definition}
Für den Übergangsvergleich sollte Symmetrie üblicherweise im normierten
Übergangsraum statt in unverarbeiteten physischen Koordinaten bewertet werden.

\subsection{Antisymmetrische Struktur}
\begin{definition}[Antisymmetrischer Übergang]
\label{def:classification-antisymmetric-transition}
Ein Übergang besitzt antisymmetrische Struktur, wenn seine Abweichung von
einem Referenz- oder Mittelwert bei Spiegelung an der Mitte ihr Vorzeichen
wechselt.
\end{definition}
Antisymmetrie ist besonders für die Analyse zentrierter
Übergangsformfunktionen oder des Ableitungsverhaltens nützlich.

\subsection{Asymmetrische Übergänge}
\begin{definition}[Asymmetrischer Übergang]
\label{def:classification-asymmetric-transition}
Ein Übergang ist asymmetrisch, wenn keine Symmetriebedingung auferlegt ist.
\end{definition}
Asymmetrie kann bewusst durch realisierungsbewussten Entwurf,
fahrzeugdynamische Optimierung, Topologie- oder lokale
Ingenieurnebenbedingungen entstehen.

\section{Ableitungsklassifikation}

\subsection{Klasse beschränkter Gradienten}
\begin{definition}[Klasse beschränkter Gradienten]
\label{def:bounded-gradient-class}
Ein Krümmungsgesetz gehört zur Klasse beschränkter Gradienten, wenn
\[
|\kappa'(s)|\le\eta_{\max}.
\]
\end{definition}

\subsection{Klasse beschränkter höherer Ableitungen}
\begin{definition}[Klasse beschränkter höherer Ableitungen]
\label{def:bounded-higher-derivative-class}
Ein Krümmungsgesetz gehört zu dieser Klasse, wenn
\[
|\kappa''(s)|<\infty,\qquad |\kappa'''(s)|<\infty.
\]
\end{definition}
Diese Klassen beeinflussen dynamische Glätte, numerische Stabilität,
Realisierungsrobustheit und Konditionierung der Optimierung.

\section{Spektrale Klassifikation}

\subsection{Niederfrequente Strukturen}
\begin{definition}[Niederfrequente Krümmungsklasse]
\label{def:low-frequency-curvature-class}
Niederfrequente Krümmungsstrukturen werden von langsam veränderlichem
Krümmungsverlauf bestimmt.
\end{definition}
Sie bestimmen üblicherweise die großräumige geometrische Form des Alignments.

\subsection{Hochfrequente Strukturen}
\begin{definition}[Hochfrequente Krümmungsklasse]
\label{def:high-frequency-curvature-class}
Hochfrequente Krümmungsstrukturen enthalten schnell veränderliche oder
oszillierende Anteile.
\end{definition}
Sie verstärken tendenziell Realisierungsempfindlichkeit, dynamische Rauheit
und Messempfindlichkeit und reduzieren die numerische Stabilität.

\begin{proposition}[Spektrales Realisierungsprinzip]
\label{prop:classification-spectral-realization}
Realisierung unterdrückt hochfrequente Krümmungsanteile.
\end{proposition}
Die spektrale Klassifikation verbindet damit Übergangstheorie und physisches
Realisierungsverhalten.

\section{Realisierungsklassifikation}

\subsection{Realisierungsstabile Klasse}
\begin{definition}[Realisierungsstabile Klasse]
\label{def:realization-stable-class}
Ein Übergang gehört zur realisierungsstabilen Klasse, wenn die dominante
Krümmungsstruktur unter Realisierung bewahrt und niederfrequentes Verhalten
identifizierbar bleibt und die Realisierung keine Instabilität einführt.
\end{definition}

\subsection{Realisierungsempfindliche Klasse}
\begin{definition}[Realisierungsempfindliche Klasse]
\label{def:realization-sensitive-class}
Ein Übergang gehört zur realisierungsempfindlichen Klasse, wenn kleine
Realisierungsstörungen sein Krümmungsverhalten stark verändern.
\end{definition}
Solche Übergänge können mathematisch gültig, aber betrieblich unerwünscht
sein.

\section{Laufzeitklassifikation}

\subsection{Projektionsstabile Klasse}
\begin{definition}[Projektionsstabile Klasse]
\label{def:projection-stable-class}
Eine Krümmungsstruktur ist projektionsstabil, wenn die
World-to-Track-Projektion numerisch robust, lokale inverse Probleme gut
konditioniert und Mehrdeutigkeiten beherrschbar bleiben.
\end{definition}

\subsection{Laufzeitglatte Klasse}
\begin{definition}[Laufzeitglatte Klasse]
\label{def:runtime-smooth-class}
Eine Krümmungsstruktur ist laufzeitglatt, wenn Frame-Bewertung stetig,
Pose-Rekonstruktion robust, Visualisierung stabil und Laufzeitoperatoren
vorhersehbar bleiben.
\end{definition}

\section{Betriebliche Klassifikation}

\subsection{Komfortorientierte Klasse}
\begin{definition}[Komfortorientierte Klasse]
\label{def:comfort-oriented-class}
Ein Übergang gehört zur komfortorientierten Klasse, wenn der Ruck klein
bleibt, sich die Querbeschleunigung glatt entwickelt und dynamische Anregung
minimiert wird.
\end{definition}

\subsection{Realisierungsorientierte Klasse}
\begin{definition}[Realisierungsorientierte Klasse]
\label{def:realization-oriented-class}
Ein Übergang gehört zur realisierungsorientierten Klasse, wenn
Baurobustheit priorisiert, Maschinenrealisierung stabil,
Instandhaltungsempfindlichkeit reduziert und das beabsichtigte Verhalten durch
physische Filterung bewahrt wird.
\end{definition}

\subsection{Optimierungsorientierte Klasse}
\begin{definition}[Optimierungsorientierte Klasse]
\label{def:optimization-oriented-class}
Ein Übergang gehört zur optimierungsorientierten Klasse, wenn die
Parametrisierung numerisch stabil, Gradienten gut konditioniert,
Nebenbedingungen ausdrückbar und die dünn besetzte Repräsentation effizient
ist.
\end{definition}

\section{Familienklassifikation}

\begin{definition}[Klassifikation von Übergangsfamilien]
\label{def:transition-family-classification}
Übergangsfamilien können nach analytischer Form, Stetigkeitsordnung,
Ableitungs- und Spektralverhalten, Realisierungsverhalten, Symmetrie,
Laufzeitstabilität und betrieblicher Eignung klassifiziert werden.
\end{definition}

Die Übergangsklassifikation ist damit mehrdimensional statt formelbasiert. Eine
Familie kann auf einer Klassifikationsachse günstig und auf einer anderen
ungünstig sein.

\section{Klassifikationsvektoren}

\begin{definition}[Klassifikationsvektor]
\label{def:classification-vector}
Einem Übergang kann ein Klassifikationsvektor zugewiesen werden:
\[
\chi(T)=(
\chi_{\mathrm{cont}},
\chi_{\mathrm{deriv}},
\chi_{\mathrm{spec}},
\chi_{\mathrm{real}},
\chi_{\mathrm{run}},
\chi_{\mathrm{op}}
),
\]
dessen Komponenten Stetigkeit, Ableitungsverhalten, Spektralstruktur,
Realisierungsverhalten, Laufzeitstabilität und betriebliche Eignung
beschreiben.
\end{definition}

Der Vektor ist keine starre Taxonomie, sondern eine strukturierte Möglichkeit,
Übergangsfamilien über mehrere relevante Dimensionen zu vergleichen.

\section{Kanonische Interpretation}

\begin{proposition}
Innerhalb von AIM ist Krümmungsklassifikation grundlegend intrinsisch und
nicht koordinatenbasiert.
\end{proposition}
\begin{proposition}
Übergangsqualität kann durch Positionsgeometrie allein nicht hinreichend
charakterisiert werden.
\end{proposition}
\begin{proposition}
Realisierung, Laufzeitstabilität, Dynamik und Optimierung erfordern eine
Klassifikation unmittelbar im Krümmungsraum.
\end{proposition}
\begin{proposition}
Die Übergangstaxonomie in AIM ist mehrdimensional und kontextabhängig.
\end{proposition}

\section{Verbindung zu AIM}

\begin{summary}
Die Krümmungsklassifikation stellt das strukturelle Interpretations-Framework
der AIM-Geometrie bereit. Übergänge werden intrinsisch klassifiziert;
Formelnamen sind gegenüber dem Krümmungsverhalten nachrangig;
Realisierungsverhalten wird analysierbar, Laufzeitrobustheit klassenabhängig,
betriebliche Eignung vergleichbar, und Optimierung wirkt unmittelbar auf
Krümmungsstrukturen.

Damit wird der Krümmungsraum zur primären Klassifikationsdomäne der
Eisenbahn-Alignment-Geometrie und zur Grundlage einer Taxonomie von
Übergangsfamilien.
\end{summary}
